\documentclass[10pt,letterpaper]{article}
\usepackage{cloudsuite}
\renewcommand{\DocShort}{CH. 13}
\begin{document}
\coverpage{Chapter 13: Fundamental Cloud Architectures}{Cornell study notes and applied review | ARCHITECTURES}{Concepts \textbullet{} decisions \textbullet{} verification artifacts}
\sourceStatement
\tableofcontents
\newpage
\section{Learning contract}
\begin{takeaway}
Select foundational architectures and their required mechanisms.
\end{takeaway}
\textbf{Chapter scope:} Distribution, pooling, elasticity, balancing, storage and multicloud.\\
\sourceRef{13}{343}
\subsection{Mapped source sections}
\begin{itemize}
\item 13.1 Workload Distribution Architecture
\item 13.2 Resource Pooling Architecture
\item 13.3 Dynamic Scalability Architecture
\item 13.4 Elastic Resource Capacity Architecture
\item 13.5 Service Load Balancing Architecture
\item 13.6 Cloud Bursting Architecture
\item 13.7 Elastic Disk Provisioning Architecture
\item 13.8 Redundant Storage Architecture
\item 13.9 Multicloud Architecture
\item 13.10 Case Study Example
\end{itemize}
\section{Cornell cue and notes matrix}
\small\setlength{\LTleft}{0pt}\setlength{\LTright}{0pt}
\begin{longtable}{>{\RaggedRight\arraybackslash}p{0.29\textwidth}>{\RaggedRight\arraybackslash}p{0.63\textwidth}}
\toprule\textbf{Cue / essential question}&\textbf{Explanatory notes and reasoning}\\\midrule\endfirsthead\toprule\textbf{Cue}&\textbf{Notes}\\\midrule\endhead
\textbf{How do pooling and distribution differ?} & Pooling makes a collective capacity resource available; distribution selects where individual requests execute. Both help utilization but solve different coordination problems.\\[3pt]\textcolor{petrol}{\emph{Practice:} Create separate block diagrams for each.}\\[8pt]\midrule
\textbf{What does dynamic scalability add?} & Scaling policies react to demand using monitors, listeners and replication. Observable signals should drive defined scaling actions.\\[3pt]\textcolor{petrol}{\emph{Practice:} State an autoscaling policy with a negative test.}\\[8pt]\midrule
\textbf{What is elastic capacity versus balancing?} & Elastic capacity changes supply, while balancing reallocates work among available instances. A service may require both.\\[3pt]\textcolor{petrol}{\emph{Practice:} Run a test with overload and a failed target.}\\[8pt]\midrule
\textbf{When would cloud bursting apply?} & Bursting extends demand handling into another environment, introducing network, synchronization and price constraints.\\[3pt]\textcolor{petrol}{\emph{Practice:} Define a maximum burst cost and data boundary.}\\[8pt]\midrule
\textbf{How is storage elasticity different from storage redundancy?} & Elastic disk provisioning changes capacity; redundant storage protects against chosen storage failures.\\[3pt]\textcolor{petrol}{\emph{Practice:} Demonstrate a capacity expansion and a disk-loss test.}\\[8pt]\midrule
\textbf{Why is multicloud not automatically resilient?} & Independent providers can still share identity, DNS, code or management dependencies.\\[3pt]\textcolor{petrol}{\emph{Practice:} Map the hidden common dependencies in a two-cloud design.}\\[8pt]\midrule
\end{longtable}\normalsize
\section{Synthesis and architecture connections}
An architecture name is a repeatable composition of mechanisms with operating assumptions. Draw component relationships, failure domains, trust boundaries and a representative workload before committing to the pattern.
\begin{workbox}[title={Applied deliverable}]
Select foundational architectures and their required mechanisms.\\[4pt]Create a one-page architecture or operations artifact with: (1) assumptions, (2) a decision and alternatives, (3) control ownership, (4) observable acceptance criteria, and (5) a failed-case test.
\end{workbox}
\section{Active recall}
\begin{enumerate}
\item How do pooling and distribution differ?
\item What does dynamic scalability add?
\item What is elastic capacity versus balancing?
\item When would cloud bursting apply?
\end{enumerate}
\subsection{Answer-check rubric}
\begin{itemize}
\item A correct response distinguishes the named concepts rather than relying on vendor names.
\item The explanation identifies an operational consequence, boundary or tradeoff.
\item At least one measurement, test or documentary artifact could demonstrate the claim.
\end{itemize}
\section{Blank Cornell capture sheet}
\textbf{Topic / source subsection:}\enspace\rule{0.65\linewidth}{0.4pt}\\[6pt]
\begin{longtable}{p{0.28\textwidth}p{0.64\textwidth}}
\toprule Cues / questions & Notes / evidence \\ \midrule
\rule{0pt}{42pt} & \\ \midrule
\rule{0pt}{42pt} & \\ \midrule
\rule{0pt}{42pt} & \\ \midrule
\rule{0pt}{42pt} & \\ \midrule
\rule{0pt}{42pt} & \\ \midrule
\end{longtable}
\textbf{Summary in three sentences:}\\[16pt]\hrulefill\\[16pt]\hrulefill\\[16pt]\hrulefill

\refs
\end{document}
